\chapter{MỞ ĐẦU}
\label{chap:introduction}

[Phần mở đầu giới thiệu bối cảnh chung của đề tài, lý do lựa chọn nghiên cứu, tính cấp thiết và ý nghĩa thực tiễn. Trích dẫn các công trình liên quan như \cite{sample_book, sample_article}.]

% ================= SECTION 1.1 =================
\section{Bối cảnh đề tài}
\label{sec:background}

[Trình bày tổng quan về lĩnh vực nghiên cứu, thực trạng vấn đề và động lực thúc đẩy thực hiện đề tài.]

% ================= SECTION 1.2 =================
\section{Mục tiêu nghiên cứu}
\label{sec:objectives}

Đề tài hướng tới việc hoàn thành các mục tiêu cụ thể sau:
\begin{itemize}
    \item \textbf{Mục tiêu 1:} [Nội dung mục tiêu nghiên cứu thứ nhất].
    \item \textbf{Mục tiêu 2:} [Nội dung mục tiêu nghiên cứu thứ hai].
    \item \textbf{Mục tiêu 3:} [Nội dung mục tiêu nghiên cứu thứ ba].
\end{itemize}

% ================= SECTION 1.3 =================
\section{Phạm vi và giới hạn đề tài}
\label{sec:scope}

\subsection{Phạm vi nghiên cứu}
[Xác định rõ giới hạn nội dung, công nghệ và đối tượng được triển khai trong phạm vi đề tài.]

\subsection{Giới hạn}
[Nêu rõ các điều kiện biên, giả định hoặc hạn chế về mặt phần cứng, môi trường thử nghiệm.]

% ================= SECTION 1.4 =================
\section{Bố cục báo cáo}
\label{sec:organization}

Báo cáo được tổ chức thành các chương như sau:
\begin{itemize}
    \item \textbf{Chương 1 -- Mở đầu:} Giới thiệu bối cảnh, mục tiêu và phạm vi nghiên cứu.
    \item \textbf{Chương 2 -- Cơ sở lý thuyết:} Trình bày các kiến thức nền tảng và mô hình toán học liên quan.
    \item \textbf{Chương 3 -- Phương pháp nghiên cứu và Kiến trúc hệ thống:} Trình bày thiết kế hệ thống và quy trình xử lý.
    \item \textbf{Chương 4 -- Hiện thực và Kết quả thực nghiệm:} Chi tiết quá trình cài đặt, ca kiểm thử và đánh giá hiệu năng.
    \item \textbf{Chương 5 -- Kết luận và Hướng phát triển:} Tổng kết các đóng góp đạt được và đề xuất hướng nghiên cứu tiếp theo.
\end{itemize}
